\documentclass{amsart}
% For dealing with references we use the comment environment
\usepackage{verbatim}
\newenvironment{reference}{\comment}{\endcomment}
%\newenvironment{reference}{}{}
\newenvironment{slogan}{\comment}{\endcomment}
\newenvironment{history}{\comment}{\endcomment}

% For commutative diagrams we use Xy-pic
\usepackage[all]{xy}

% We use 2cell for 2-commutative diagrams.
\xyoption{2cell}
\UseAllTwocells

% We use multicol for the list of chapters between chapters
\usepackage{multicol}

% This is generally recommended for better output
\usepackage{lmodern}
\usepackage[T1]{fontenc}

% For cross-file-references

% Package for hypertext links:
\usepackage{hyperref}

% For any local file, say "hello.tex" you want to link to please

% Theorem environments.
%
\theoremstyle{plain}
\newtheorem{theorem}[subsection]{Theorem}
\newtheorem{proposition}[subsection]{Proposition}
\newtheorem{lemma}[subsection]{Lemma}

\theoremstyle{definition}
\newtheorem{definition}[subsection]{Definition}
\newtheorem{example}[subsection]{Example}
\newtheorem{exercise}[subsection]{Exercise}
\newtheorem{situation}[subsection]{Situation}

\theoremstyle{remark}
\newtheorem{remark}[subsection]{Remark}
\newtheorem{remarks}[subsection]{Remarks}

\numberwithin{equation}{subsection}

% Macros
%
\def\lim{\mathop{\mathrm{lim}}\nolimits}
\def\colim{\mathop{\mathrm{colim}}\nolimits}
\def\Spec{\mathop{\mathrm{Spec}}}
\def\Hom{\mathop{\mathrm{Hom}}\nolimits}
\def\Ext{\mathop{\mathrm{Ext}}\nolimits}
\def\SheafHom{\mathop{\mathcal{H}\!\mathit{om}}\nolimits}
\def\SheafExt{\mathop{\mathcal{E}\!\mathit{xt}}\nolimits}
\def\Sch{\mathit{Sch}}
\def\Mor{\mathop{\mathrm{Mor}}\nolimits}
\def\Ob{\mathop{\mathrm{Ob}}\nolimits}
\def\Sh{\mathop{\mathit{Sh}}\nolimits}
\def\NL{\mathop{N\!L}\nolimits}
\def\CH{\mathop{\mathrm{CH}}\nolimits}
\def\proetale{{pro\text{-}\acute{e}tale}}
\def\etale{{\acute{e}tale}}
\def\QCoh{\mathit{QCoh}}
\def\Ker{\mathop{\mathrm{Ker}}}
\def\Im{\mathop{\mathrm{Im}}}
\def\Coker{\mathop{\mathrm{Coker}}}
\def\Coim{\mathop{\mathrm{Coim}}}

% Boxtimes
%
\DeclareMathSymbol{\boxtimes}{\mathbin}{AMSa}{"02}

%
% Macros for moduli stacks/spaces
%
\def\QCohstack{\mathcal{QC}\!\mathit{oh}}
\def\Cohstack{\mathcal{C}\!\mathit{oh}}
\def\Spacesstack{\mathcal{S}\!\mathit{paces}}
\def\Quotfunctor{\mathrm{Quot}}
\def\Hilbfunctor{\mathrm{Hilb}}
\def\Curvesstack{\mathcal{C}\!\mathit{urves}}
\def\Polarizedstack{\mathcal{P}\!\mathit{olarized}}
\def\Complexesstack{\mathcal{C}\!\mathit{omplexes}}
% \Pic is the operator that assigns to X its picard group, usage \Pic(X)
% \Picardstack_{X/B} denotes the Picard stack of X over B
% \Picardfunctor_{X/B} denotes the Picard functor of X over B
\def\Pic{\mathop{\mathrm{Pic}}\nolimits}
\def\Picardstack{\mathcal{P}\!\mathit{ic}}
\def\Picardfunctor{\mathrm{Pic}}
\def\Deformationcategory{\mathcal{D}\!\mathit{ef}}

\setlength{\emergencystretch}{3em}
\begin{document}
\title[Stacks Project, Tag 0DYG]{344 --- Testing formal smoothness by small Artinian extensions}
\maketitle
\noindent Copyright (C) 2005--2025 Johan de Jong. Stacks Project authors. Source: \href{https://stacks.math.columbia.edu/tag/0DYG}{Tag 0DYG}.

\noindent Editorial difficulty note (not author text). Difficulty H2: The proof reduces arbitrary lifting diagrams to local Artinian ones with the same residue field. For square-zero kernels of arbitrary dimension it forms a product of one-dimensional quotient extensions and splits the resulting vector-space kernel to recover a lift..

\noindent Editorial provenance note (not author text). History: excerpt from revision \texttt{a0444\allowbreak 6e57e\allowbreak c1fbc\allowbreak 252a8\allowbreak 71afc\allowbreak ec775\allowbreak 2fb28\allowbreak 07b14}, more-algebra.tex, lines 9435--9528. Reference presentation was adapted; original-author context and core support blocks were retained or added. The mathematical wording is unchanged. Licensed under GNU FDL 1.2 or later; no invariant sections or cover texts. See ../licenses/COPYING. Context blocks reproduce author definitions and supporting results; they are not additional counted problems.

\begin{definition}
\label{algebra-definition-small-extension}
Let $\varphi : B' \to B$ be a ring map.
We say $\varphi$ is a {\it small extension} if
$B'$ and $B$ are local Artinian rings, $\varphi$ is surjective
and $I = \Ker(\varphi)$ has length $1$ as a $B'$-module.
\end{definition}

\begin{definition}
\label{definition-formally-smooth}
Let $R \to S$ be a homomorphism of topological rings with $R$ and $S$
linearly topologized. We say $S$ is {\it formally smooth over $R$} if
for every commutative solid diagram
$$
\xymatrix{
S \ar[r] \ar@{-->}[rd] & A/J \\
R \ar[r] \ar[u] & A \ar[u]
}
$$
of homomorphisms of topological rings where $A$ is a discrete ring and
$J \subset A$ is an ideal of square zero, a dotted arrow exists which
makes the diagram commute.
\end{definition}

\begin{lemma}
\label{lemma-fs-local}
Let $(R, \mathfrak m) \to (S, \mathfrak n)$ be a local homomorphism
of local rings. The following are equivalent
\begin{enumerate}
\item $R \to S$ is formally smooth in the $\mathfrak n$-adic topology,
\item for every solid commutative diagram
$$
\xymatrix{
S \ar[r] \ar@{-->}[rd] & A/J \\
R \ar[r] \ar[u] & A \ar[u]
}
$$
of local homomorphisms of local rings where $J \subset A$ is
an ideal of square zero, $\mathfrak m_A^n = 0$ for some $n > 0$, and
$S \to A/J$ induces an isomorphism on residue fields, a dotted
arrow exists which makes the diagram commute.
\end{enumerate}
If $S$ is Noetherian these conditions are also equivalent to
\begin{enumerate}
\item[(3)] same as in (2) but only for diagrams where in addition
$A \to A/J$ is a small extension
(Algebra, Definition \ref{algebra-definition-small-extension}).
\end{enumerate}
\end{lemma}

\begin{proof}
The implication (1) $\Rightarrow$ (2) follows from the definitions.
Consider a diagram
$$
\xymatrix{
S \ar[r] \ar@{-->}[rd] & A/J \\
R \ar[r] \ar[u] & A \ar[u]
}
$$
as in Definition \ref{definition-formally-smooth} for the
$\mathfrak m$-adic topology on $R$ and the $\mathfrak n$-adic topology on $S$.
Pick $m > 0$ with $\mathfrak n^m(A/J) = 0$
(possible by continuity of maps in diagram).
Consider the subring $A'$ of $A$ which is the inverse image
of the image of $S$ in $A/J$. Set $J' = J$ viewed as an ideal in $A'$.
Then $J'$ is an ideal of square zero in $A'$ and $A'/J'$
is a quotient of $S/\mathfrak n^m$. Hence $A'$ is local
and $\mathfrak m_{A'}^{2m} = 0$. Thus we get a diagram
$$
\xymatrix{
S \ar[r] \ar@{-->}[rd] & A'/J' \\
R \ar[r] \ar[u] & A' \ar[u]
}
$$
as in (2). If we can construct the dotted arrow in this diagram,
then we obtain the dotted arrow in the original one by composing
with $A' \to A$. In this way we see that (2) implies (1).

\medskip\noindent
Assume $S$ Noetherian. The implication (1) $\Rightarrow$ (3) is immediate.
Assume (3) and suppose a diagram as in (2) is given.
Then $\mathfrak m_A^n J = 0$ for some $n > 0$.
Considering the maps
$$
A \to A/\mathfrak m_A^{n - 1}J \to \ldots \to A/\mathfrak mJ \to A/J
$$
we see that it suffices to produce the lifting if $\mathfrak m_A J = 0$.
Assume $\mathfrak m_A J = 0$ and let $A' \subset A$ be the ring
constructed above. Then $A'/J'$ is Artinian
as a quotient of the Artinian local ring $S/\mathfrak n^m$.
Thus it suffices to show that given property (3) we can find the dotted
arrow in diagrams as in (2) with $A/J$ Artinian and
$\mathfrak m_A J = 0$.
Let $\kappa$ be the common residue field of $A$,
$A/J$, and $S$. By (3), if $J_0 \subset J$ is an ideal with
$\dim_\kappa(J/J_0) = 1$, then we can produce a dotted arrow
$S \to A/J_0$. Taking the product we obtain
$$
S \longrightarrow \prod\nolimits_{J_0 \text{ as above}} A/J_0
$$
Clearly the image of this arrow is contained in the sub $R$-algebra
$A'$ of elements which map into the small diagonal
$A/J \subset \prod_{J_0} A/J$.
Let $J' \subset A'$ be the elements mapping to zero in $A/J$.
Then $J'$ is an ideal of square zero and as $\kappa$-vector space
equal to
$$
J' = \prod\nolimits_{J_0 \text{ as above}} J/J_0
$$
Thus the map $J \to J'$ is injective. By the theory of vector spaces
we can choose a splitting $J' = J \oplus M$. It follows that
$$
A' = A \oplus M
$$
as an $R$-algebra. Hence the map $S \to A'$ can be composed with
the projection $A' \to A$ to give the desired dotted arrow thereby
finishing the proof of the lemma.
\end{proof}
\end{document}
